% bank_fintech_locus/drafts/zh/manuscript.tex
% P4 中文 manuscript 骨架 — M0 G0 scaffolded per B-route plan v2.55 (2026-05-17)
% 编译: xelatex + biber + xelatex × 2

% papers_assets/_shared/preamble_zh.tex
% Shared Chinese LaTeX preamble (XeLaTeX + CJKutf8).
% P4.1 (plan-r3 derivative-paper extension, 2026-05-15)
% Compile: xelatex + biber + xelatex × 2

\documentclass[12pt,a4paper]{article}

% --- CJK Fonts (XeLaTeX, cross-platform Q1.1b fallback) ---
\usepackage{fontspec}
\usepackage{xeCJK}
% Windows (local dev): SimSun/SimHei/FangSong
% Linux/macOS (CI): Noto Serif CJK SC / Noto Sans CJK SC (fonts-noto-cjk apt package)
% Last resort: AR PL UMing CN (fonts-arphic-uming) which is preinstalled in many texlive distros
\IfFontExistsTF{SimSun}{%
  \setCJKmainfont{SimSun}[BoldFont=SimHei]%
}{%
  \IfFontExistsTF{Noto Serif CJK SC}{%
    \setCJKmainfont{Noto Serif CJK SC}[BoldFont=Noto Serif CJK SC]%
  }{%
    \setCJKmainfont{AR PL UMing CN}[AutoFakeBold=2]%
  }%
}
\IfFontExistsTF{SimHei}{%
  \setCJKsansfont{SimHei}%
}{%
  \IfFontExistsTF{Noto Sans CJK SC}{%
    \setCJKsansfont{Noto Sans CJK SC}%
  }{%
    \setCJKsansfont{AR PL UKai CN}%
  }%
}
\IfFontExistsTF{FangSong}{%
  \setCJKmonofont{FangSong}%
}{%
  \IfFontExistsTF{Noto Sans Mono CJK SC}{%
    \setCJKmonofont{Noto Sans Mono CJK SC}%
  }{%
    \setCJKmonofont{AR PL UMing CN}%
  }%
}
\setmainfont{Times New Roman}[Ligatures=TeX]
\xeCJKsetup{CJKecglue={\hskip 0.15em plus 0.05em minus 0.05em}}

% --- Math ---
\usepackage{amsmath, amssymb, amsthm}

% --- Tables & Figures ---
\usepackage{booktabs}
\usepackage{graphicx}
\usepackage{longtable}
\usepackage{rotating}
\usepackage{caption}
\captionsetup{font=small,labelfont=bf,skip=4pt}

% --- Layout ---
\usepackage[margin=1in]{geometry}
\usepackage{setspace}
\onehalfspacing
\usepackage{hyperref}
\hypersetup{colorlinks=true,linkcolor=blue,citecolor=blue,urlcolor=blue}

% --- Bibliography ---
\usepackage[backend=biber,style=apa,natbib=true]{biblatex}

% --- Macros ---
\newcommand{\code}[1]{\texttt{#1}}
\usepackage[dvipsnames]{xcolor}
\newcommand{\NoteThis}[1]{\textbf{\textcolor{red}{[TODO: #1]}}}

% --- Theorem environments ---
\newtheorem{prop}{命题}
\newtheorem{hyp}{假设}


\addbibresource{refs.bib}

\title{银行 fintech 公告之 Locus：\\三波浪潮下底层技术 vs 应用场景的市场反应差异}
\author{\NoteThis{作者姓名}}
\date{\today}

\begin{document}
\maketitle

\begin{abstract}
\NoteThis{最终摘要待定 — 目标约 300 中文字。}
本文遵循 Gomber 等 (2018, \emph{JMIS}) 的 fintech 分类学, 把银行 fintech 公告区分为
底层技术 (区块链 / AI / 云) 与应用场景 (支付 / 信贷 / 理财) 两类 locus 属性。基于
2007--2024 (18 年) 的中国银行业面板, 检验此二元 locus 如何影响事件窗口 5 日 CAR、
专利产出 (Kogan 等 2017 市值加权) 及 IT capex / R\&D 中介通道。识别策略利用 4 个外生
政策冲击 (余额宝、蚂蚁 IPO 暂停、两项数据出境管理办法) 与三波 fintech 浪潮 (互联网
金融 / 区块链 / 大模型, 以 consolidation period 作对照)。研究发现 \NoteThis{核心结论}\@。
\end{abstract}

\section{引言}\label{sec:intro}
\NoteThis{研究问题动机与贡献。}

\section{文献综述}\label{sec:lit}
\NoteThis{三条文献: (1) Fintech 分类与 locus \autocite{gomber2018fintech,goldstein2019fintech};
(2) 创新市值加权 \autocite{kogan2017technological};
(3) 区块链与代币化 \autocite{conghe2019blockchain,cong2021tokenomics}。}

\section{方法}\label{sec:method}
\subsection{样本与数据}
$N \times T \approx 3000\text{--}9000$ 银行—年, 2007--2024。

\subsection{Locus 度量 (stage 32a)}
\NoteThis{innovation\_locus\_full 二元指示 → derived primary\_pair
(locus\_layer\_full, locus\_scenario\_full) 与 locus\_gap.}

\subsection{识别策略}
\NoteThis{Layer $\times$ event study + leave-one-out peer-locus 工具变量;
4 events $\times$ 5 日 CAR + 专利 + Kogan 2017 市值加权。}

\subsection{嵌入类测度 (双主轨、互为稳健性; F5)}
所有嵌入派生特征均在来自独立厂商的两个同为主测度的嵌入模型下并行
计算、并排报告、互为稳健性: 阿里 \code{Qwen3-Embedding-8B}
(4{,}096 维, 经硅基流动 API; canonical 无后缀列) 与 OpenAI
\code{text-embedding-3-large} (3{,}072 维; \code{\_lg} 后缀列),
遵循严格列对偶 (任一特征两轨齐备或皆无)。各轨进入各自恰好识别的
独立设定, 绝不合并进同一过度识别系统。早期管线版本 canonical 槽为
Google Gemini 嵌入; 自管线 v3.15.0 (2026 年 6 月) 起换为阿里
\code{Qwen3-Embedding-8B}, 全部嵌入派生特征经受控全量重嵌再生,
嵌入空间语义效度以 Spearman 秩相关闸把关 ($\rho \geq 0.80$)。
情感标注 LLM 集成不受此变更影响: 嵌入模型与情感 LLM
是互不相交的管线组件。

\section{数据}\label{sec:data}
\NoteThis{数据来源 (CSMAR 日行情, IncoPat 专利, Kogan 2017 算法); 描述统计 (Table 1)。}

\section{实证结果}\label{sec:results}
\NoteThis{Table 2 (主回归: layer $\times$ event study); Table 3 (Baron-Kenny 中介:
locus $\to$ R\&D / IT capex $\to$ patent + CAR); Table 4 (异质性 三波 fintech 浪潮
wave\_1/2/3 + consolidation + SOE); Table 5 (稳健性: alternative locus classifier +
Borusyak DiD)。}

\section{结论}\label{sec:conclusion}
\NoteThis{总结发现; fintech 监管与银行创新战略启示。}

\printbibliography

\end{document}
